%%=============================================================
%% chapter5_4.tex
%% Phần 5.4 — Kết quả đánh giá: AI #3 Phát hiện vật thể
%%
%% \input file — cần thêm vào preamble:
%%   \usepackage{tikz}
%%   \usetikzlibrary{arrows.meta,positioning,patterns}
%%   \usepackage{pgfplots}
%%   \pgfplotsset{compat=1.18}
%%   \usepackage{booktabs,multirow,subcaption,xcolor}
%%=============================================================

% ================================================================
\section{Kết quả và Đánh giá: AI \#3 --- Phát hiện vật thể}
\label{sec:ai3-results}
% ================================================================

% ----------------------------------------------------------------
\subsection{Kết quả detection model: COCO metrics}
\label{subsec:ai3-coco}
% ----------------------------------------------------------------

Mô hình được đánh giá trên cả validation set và test set độc lập
theo chuẩn \textbf{COCO evaluation protocol} --- tính trung bình mAP qua
các ngưỡng IoU từ 0.50 đến 0.95 với bước 0.05.

\begin{figure}[htbp]
  \centering
  \begin{tikzpicture}
    \begin{axis}[
      ybar, bar width=0.45cm,
      width=10cm, height=7cm,
      symbolic x coords={mAP (0.50:0.95), AP50, AP75, AR},
      xtick=data,
      x tick label style={font=\small, align=center},
      ymin=0, ymax=1.15,
      ylabel={Score},
      ylabel style={font=\small},
      legend style={at={(0.5,1.03)}, anchor=south, legend columns=2, font=\small},
      legend cell align=left,
      nodes near coords,
      every node near coord/.style={font=\scriptsize, /pgf/number format/fixed,
                                    /pgf/number format/precision=3},
      enlarge x limits=0.25,
      grid=major, grid style={dashed, gray!30},
      axis line style={gray!60},
      tick style={gray!60},
    ]
    \addplot[fill=cppblue!70, draw=cppblue!90] coordinates {
      (mAP (0.50:0.95), 0.791)
      (AP50, 0.990)
      (AP75, 0.990)
      (AR,   0.869)
    };
    \addplot[fill=orng!80, draw=orng!90] coordinates {
      (mAP (0.50:0.95), 0.799)
      (AP50, 1.000)
      (AP75, 0.969)
      (AR,   0.856)
    };
    \legend{Validation, Test}
    \end{axis}
  \end{tikzpicture}
  \caption{COCO metrics trên validation set và test set. AP50 trên test set
           đạt 1.000 cho thấy 100\% cylinder được phát hiện đúng ở ngưỡng IoU
           lỏng lẻ (50\%). mAP$\approx$0.80 phản ánh vẫn còn đây đó những trường hợp
           định vị bounding box chưa được chính xác ở ngưỡng IoU cao.}
  \label{fig:ai3-coco}
\end{figure}

Bảng~\ref{tab:ai3-metrics} tổng hợp số liệu chi tiết:

\begin{table}[htbp]
  \centering
  \caption{COCO metrics trên validation và test set}
  \label{tab:ai3-metrics}
  \small
  \begin{tabular}{lcc}
    \toprule
    \textbf{Metric}                    & \textbf{Validation} & \textbf{Test} \\
    \midrule
    mAP (IoU 0.50:0.95)               & 0.791               & \textbf{0.799} \\
    AP50 (IoU $\geq$ 0.50)            & 0.990               & \textbf{1.000} \\
    AP75 (IoU $\geq$ 0.75)            & 0.990               & 0.969 \\
    AR (maxDets=100)                   & 0.869               & 0.856 \\
    \midrule
    Precision (IoU=0.50, conf=0.50)    & ---                 & \textbf{1.000} \\
    Recall    (IoU=0.50, conf=0.50)    & ---                 & 0.990 \\
    F1-Score                           & ---                 & \textbf{0.995} \\
    TP / FP / FN                       & ---                 & 95 / 0 / 1 \\
    \bottomrule
  \end{tabular}
\end{table}

\begin{figure}[htbp]
  \centering
  \begin{tikzpicture}
    \begin{axis}[
      ybar, bar width=0.55cm,
      width=8.5cm, height=6cm,
      symbolic x coords={Precision, Recall, F1-Score},
      xtick=data,
      x tick label style={font=\small},
      ymin=0.96, ymax=1.04,
      ylabel={Score},
      ylabel style={font=\small},
      nodes near coords,
      every node near coord/.style={font=\small, /pgf/number format/fixed,
                                    /pgf/number format/precision=3},
      enlarge x limits=0.3,
      grid=major, grid style={dashed, gray!30},
      axis line style={gray!60},
    ]
    \addplot[fill=pygreen!65, draw=pygreen!80] coordinates {
      (Precision, 1.000)
      (Recall,    0.990)
      (F1-Score,  0.995)
    };
    \end{axis}
  \end{tikzpicture}
  \caption{Precision, Recall và F1-Score trên test set (PASCAL VOC protocol,
           IoU$\geq$0.50, confidence$\geq$0.50, chỉ lớp cylinder).
           Precision = 1.000 và FP = 0 chứng tỏ không có false positive
           nào ở ngưỡng 0.50. FN = 1 (một cylinder bị bỏ sót) giải thích
           Recall = 0.990.}
  \label{fig:ai3-prf}
\end{figure}

Giá trị AP50 = 1.000 trên test set à ngưỡng IoU lỏng lẻ (50\%) phản ánh
mô hình phát hiện được toàn bộ cylinder. mAP = 0.799 (trung bình
qua IoU 0.50--0.95) thấp hơn một chút do mô hình EfficientDet-Lite0
đơn giản vẫn còn giới hạn trong việc định vị bounding box chính xác
ở các ngưỡng IoU cao ($\geq$0.75) --- điều phổ biến với các kiến trúc
nhe hơn.


% ----------------------------------------------------------------
\subsection{Phân tích confidence score: true positive vs false positive}
\label{subsec:ai3-conf}
% ----------------------------------------------------------------

Một trong những điều quan trọng khi triển khai thực tế là khoảng cách
phân tách giữa true positive và false positive --- điều này quyết định
ngưỡng confidence nào thực tế dùng được mà không sai kết quả.

Kiểm tra trên 218 ảnh cylinder (true positives) và 10 ảnh cube
(false positives):

\begin{figure}[htbp]
  \centering
  \begin{tikzpicture}
    \begin{axis}[
      width=11cm, height=5.5cm,
      xlabel={Confidence score},
      ylabel={},
      xmin=0, xmax=1.05,
      ymin=0, ymax=3,
      xtick={0,0.1,...,1.0},
      ytick={0.7, 2.0},
      yticklabels={\small False Positive\\(cube, $n$=10), \small True Positive\\(cylinder, $n$=218)},
      y tick label style={align=right, font=\small},
      x tick label style={font=\small},
      xlabel style={font=\small},
      grid=both, grid style={dashed, gray!25},
      axis line style={gray!60},
    ]

    %% --- True Positive band (y=2) ---
    %% Most scores: 0.85-0.98 (majority)
    \addplot[only marks, mark=|, mark size=5pt, mark options={draw=pygreen!70, line width=1.2pt},
             opacity=0.5]
    coordinates {
      (0.29,2)(0.61,2)(0.74,2)
      (0.82,2)(0.83,2)(0.84,2)(0.85,2)(0.86,2)(0.87,2)(0.88,2)
      (0.89,2)(0.90,2)(0.91,2)(0.92,2)(0.93,2)(0.94,2)(0.94,2)
      (0.95,2)(0.95,2)(0.96,2)(0.96,2)(0.97,2)(0.97,2)(0.98,2)
    };

    %% --- False Positive band (y=0.7) ---
    \addplot[only marks, mark=|, mark size=5pt, mark options={draw=warnred!70, line width=1.2pt},
             opacity=0.6]
    coordinates {
      (0.08,0.7)(0.13,0.7)(0.17,0.7)(0.21,0.7)
      (0.25,0.7)(0.28,0.7)(0.30,0.7)(0.33,0.7)(0.35,0.7)(0.36,0.7)
    };

    %% --- Threshold line at 0.50 ---
    \addplot[thick, dashed, draw=gray!80, domain=0:3]
    coordinates {(0.50, 0) (0.50, 3)};
    \node[font=\scriptsize, text=gray!80, rotate=90]
         at (axis cs: 0.52, 1.5) {threshold = 0.50};

    %% --- Gap annotation ---
    \draw[->, >=Stealth, thick, gray!70]
         (axis cs:0.38, 2.5) -- (axis cs:0.48, 2.5);
    \draw[->, >=Stealth, thick, gray!70]
         (axis cs:0.34, 2.5) -- (axis cs:0.27, 2.5);
    \node[font=\scriptsize, text=gray!80] at (axis cs:0.37, 2.7) {gap = 0.34};

    %% --- Median annotation for TP ---
    \addplot[only marks, mark=diamond*, mark size=4pt,
             mark options={draw=pygreen!80!black, fill=pygreen!60}]
    coordinates {(0.94, 2)};
    \node[font=\scriptsize, text=pygreen!80!black, above]
         at (axis cs:0.94, 2) {median 0.94};

    %% --- Max FP annotation ---
    \addplot[only marks, mark=diamond*, mark size=4pt,
             mark options={draw=warnred!80, fill=warnred!50}]
    coordinates {(0.36, 0.7)};
    \node[font=\scriptsize, text=warnred, below]
         at (axis cs:0.36, 0.7) {max 0.36};

    \end{axis}
  \end{tikzpicture}
  \caption{Phân tách confidence score giữa true positive (cylinder, xanh lá)
           và false positive (cube, đỏ). Mỗi vạch đứng đại diện một detection.
           Ngưỡng 0.50 nằm chính giữa khoảng gap 0.34,
           loại bỏ toàn bộ false positive mà vẫn chấp nhận 99.5\% true positive.}
  \label{fig:ai3-conf}
\end{figure}

Khoảng cách 0.34 giữa điểm false positive cao nhất (0.36) và 99.1\%
true positive ($\geq$0.70) xác nhận mô hình đã học được đặc trưng
phân biệt mạnh giữa cylinder và cube. Ngưỡng hiện tại 0.50 nằm an toàn
giữa hai vùng phân bố này.

Ba ảnh outlier có score thấp (0.29, 0.61, 0.74) tương ứng với cylinder
bị nhìn từ góc bất thường hoặc bị che khuất một phần. Ỡ runtime,
các trường hợp này fallback về click thủ công --- hành vi an toàn.

\begin{table}[htbp]
  \centering
  \caption{Phân tích confidence score trên 218 ảnh cylinder + 10 ảnh cube}
  \label{tab:ai3-conf}
  \small
  \begin{tabular}{lcc}
    \toprule
    & \textbf{True Positive (cylinder)} & \textbf{False Positive (cube)} \\
    \midrule
    Score range      & 0.29 -- 0.98          & 0.08 -- 0.36 \\
    Median           & \textbf{0.94}         & --- \\
    $\geq$ 0.50      & 99.5\% (217/218)      & \textbf{0/10 (0\%)} \\
    $\geq$ 0.70      & 99.1\% (216/218)      & 0/10 (0\%) \\
    $\geq$ 0.80      & 95.9\% (209/218)      & 0/10 (0\%) \\
    \midrule
    Gap ph\^an t\'ach & \multicolumn{2}{c}{$0.70 - 0.36 = 0.34$} \\
    \bottomrule
  \end{tabular}
\end{table}


% ----------------------------------------------------------------
\subsection{Kết quả tích hợp hệ thống}
\label{subsec:ai3-integration}
% ----------------------------------------------------------------

\begin{table}[htbp]
  \centering
  \caption{So sánh khả năng trước và sau khi tích hợp AI}
  \label{tab:ai3-comparison}
  \small
  \begin{tabular}{lll}
    \toprule
    \textbf{Tiêu chí}                    & \textbf{Hệ thống gốc} & \textbf{Sau cải tiến} \\
    \midrule
    Khởi tạo mỗi chu kỳ               & Click thủ công (bắt buộc) & Tự động (AI) \\
    Tỷ lệ init thành công             & 100\% (do manual)  & $\sim$99.5\% (AI) \\
    Fallback khi AI thất bại           & Không có              & Click thủ công \\
    Phục hồi khi mất tracking          & Reset toàn bộ        & AI re-detect tự động \\
    Chạy không giám sát                & Không thể              & Có (trường hợp bình thường) \\
    \bottomrule
  \end{tabular}
\end{table}

Thiết kế hiện tại gọi AI trực tiếp \textbf{mỗi iteration} của servo
loop thay vì dùng blob tracker. Điều này có một hàm ý đáng chú ý:
\textit{không có tracking failure} theo nghĩa truyền thống --- mỗi frame
là một lần detection độc lập. Tính ổn định của servo phụ thuộc vào
sự nhất quán của AI detection qua các frame liên tiếp. Cơ chế \texttt{--hint}
(truyền vị trí frame trước) giúp được phần nào khi có nhiều cylinder
trong cảnh.


% ----------------------------------------------------------------
\subsection{Phân tích latency: subprocess CPU vs EdgeTPU}
\label{subsec:ai3-latency}
% ----------------------------------------------------------------

Mỗi lần gọi \texttt{detectCylinderWithAI()} bao gồm nhiều thành phần
latency khác nhau. Điều ít người để ý là: phần overhead lớn nhất
không phải là inference mà là \textbf{Python startup time} --- chi phí
mở một process mới mỗi lần gọi \texttt{popen()}.

\begin{figure}[htbp]
  \centering
  \begin{tikzpicture}
    \begin{axis}[
      xbar stacked,
      bar width=0.55cm,
      width=11cm, height=6.5cm,
      symbolic y coords={
        {EdgeTPU\\(persistent, tương lai)},
        {CPU\\(persistent, tương lai)},
        {EdgeTPU\\(subprocess hiện tại)},
        {CPU\\(subprocess hiện tại)}
      },
      ytick=data,
      y tick label style={align=right, font=\small},
      xlabel={Thời gian (ms)},
      xlabel style={font=\small},
      xmin=0, xmax=780,
      legend style={at={(0.5,-0.18)}, anchor=north, legend columns=4,
                    font=\scriptsize},
      legend cell align=left,
      grid=major, grid style={dashed, gray!25},
      axis line style={gray!60},
      enlarge y limits=0.18,
      nodes near coords, nodes near coords align={horizontal},
      point meta=explicit symbolic,
      every node near coord/.style={font=\scriptsize},
    ]
    %%  stacks: I/O | Inference | Model load | Python startup
    %% CPU subprocess: 10 + 150 + 100 + 350 = 610
    \addplot[fill=gray!40, draw=gray!60] coordinates {
      (10, {CPU\\(subprocess hiện tại)})
      (10, {EdgeTPU\\(subprocess hiện tại)})
      (10, {CPU\\(persistent, tương lai)})
      (5,  {EdgeTPU\\(persistent, tương lai)})
    };  %% I/O
    \addplot[fill=warnred!55, draw=warnred!70] coordinates {
      (150, {CPU\\(subprocess hiện tại)})
      (10,  {EdgeTPU\\(subprocess hiện tại)})
      (150, {CPU\\(persistent, tương lai)})
      (10,  {EdgeTPU\\(persistent, tương lai)})
    };  %% Inference
    \addplot[fill=orng!65, draw=orng!80] coordinates {
      (100, {CPU\\(subprocess hiện tại)})
      (50,  {EdgeTPU\\(subprocess hiện tại)})
      (0,   {CPU\\(persistent, tương lai)})
      (0,   {EdgeTPU\\(persistent, tương lai)})
    };  %% Model load
    \addplot[fill=cppblue!55, draw=cppblue!70] coordinates {
      (350, {CPU\\(subprocess hiện tại)})
      (350, {EdgeTPU\\(subprocess hiện tại)})
      (0,   {CPU\\(persistent, tương lai)})
      (0,   {EdgeTPU\\(persistent, tương lai)})
    };  %% Python startup
    \legend{I/O (JPEG), Inference, Model load, Python startup}
    \end{axis}
  \end{tikzpicture}
  \caption{Phân tích latency mỗi lần gọi AI detection.
           Python startup time ($\sim$350~ms) chiếm phần lẫn tổng overhead
           trong kiến trúc subprocess hiện tại --- lớn hơn nhiều so với
           thời gian inference thực tế. Kiến trúc persistent process (tương lai)
           loại bỏ hoàn toàn chi phí này.}
  \label{fig:ai3-latency}
\end{figure}

\begin{table}[htbp]
  \centering
  \caption{So sánh tổng latency giữa các cấu hình triển khai}
  \label{tab:ai3-latency}
  \small
  \begin{tabular}{lcc}
    \toprule
    \textbf{Cấu hình}                    & \textbf{Inference} & \textbf{Tổng / lần gọi} \\
    \midrule
    CPU subprocess (hiện tại)            & $\sim$100--200~ms  & $\sim$600--700~ms \\
    EdgeTPU subprocess (hiện tại)       & $\sim$4--15~ms     & $\sim$420--430~ms \\
    CPU persistent (tương lai)           & $\sim$100--200~ms  & $\sim$160--210~ms \\
    EdgeTPU persistent (tương lai)      & $\sim$4--15~ms     & $\sim$15--20~ms  \\
    \bottomrule
  \end{tabular}
\end{table}

Kết quả trên cho thấy rõ điểm nghẻn hiện tại:
EdgeTPU inference chỉ mất 4--15~ms nhưng toàn bộ lá gọi
vẫn tốn 420~ms do phải khởi động Python mỗi lần. Với
kiến trúc \textbf{persistent Python process} --- một process
duy nhất khởi động một lần, load mô hình vào bộ nhớ, sau
đó tiếp nhận request từ C++ qua stdin/stdout pipe --- latency
mỗi lần gọi giảm xuống còn $\sim$15--20~ms với EdgeTPU,
mở ra khả năng chạy AI detection \textit{mỗi frame} camera
thay vì chỉ khi cần.

% Chú thích về nguồn số liệu
\smallskip
\noindent\textit{\small
Lưu ý: COCO metrics trong Bảng~\ref{tab:ai3-metrics} và Hình~\ref{fig:ai3-coco}
lấy từ kết quả đánh giá hai lớp (cylinder + cube) trên repository huấn luyện.
Phân tích confidence score trong Bảng~\ref{tab:ai3-conf} lấy từ kiểm tra
độc lập trên 218 ảnh cylinder và 10 ảnh cube (CPU inference via ai-edge-litert).
}
